\frfilename{mt1a2.tex}
\versiondate{21.11.03}
\copyrightdate{1994}

\def\chaptername{Lampiran}
\def\sectionname{Himpunan terbuka dan tertutup dalam $\BbbR^r$}

\newsection{1A2}

\def\headlinesectionname{Himpunan terbuka dan tertutup dalam $\eightBbb R^r$}

Dalam 111G saya memberikan definisi himpunan terbuka dalam $\Bbb R$ atau
$\BbbR^r$, dan dalam 121D saya menggunakan secara sepintas beberapa sifat
dasar himpunan-himpunan ini; barangkali akan membantu jika saya memaparkan
sedikit saja teori elementernya.

\leader{1A2A}{Himpunan terbuka} Ingat bahwa suatu himpunan
$G\subseteq\Bbb R$ disebut {\bf terbuka} jika untuk setiap $x\in G$ terdapat
suatu $\delta>0$ sedemikian sehingga
$\ooint{x-\delta,x+\delta}\subseteq G$; demikian pula, suatu himpunan
$G\subseteq\BbbR^r$ disebut {\bf terbuka} jika untuk setiap $x\in G$ terdapat
suatu $\delta>0$ sedemikian sehingga $U(x,\delta)\subseteq G$, dengan
$U(x,\delta)=\{y:\|y-x\|<\delta\}$, dan kita menulis $\|z\|$ untuk
$\sqrt{\zeta_1^2+\ldots+\zeta_r^2}$ jika
$z=(\zeta_1,\ldots,\zeta_r)$. \cmmnt{Selanjutnya saya memberikan argumen
untuk $r$ umum; jika saat ini Anda hanya tertarik pada kasus satu dimensi,
Anda seharusnya tidak mengalami kesulitan membacanya seolah-olah $r=1$
sepanjang uraian.}

\vleader{36pt}{1A2B}{Keluarga semua himpunan terbuka}
Misalkan $\frak T$ adalah keluarga himpunan-himpunan terbuka dalam $\BbbR^r$.
\cmmnt{Maka $\frak T$ mempunyai sifat-sifat berikut.}

(a) $\emptyset\in\frak T$\cmmnt{, yaitu himpunan kosong terbuka}.
\prooflet{\Prf\ Karena definisi `$\emptyset$ terbuka' dimulai dengan `untuk
setiap $x\in\emptyset$, $\ldots$', definisi itu pasti terpenuhi secara vakum
untuk himpunan kosong.\ \Qed}

(b) $\BbbR^r\in\frak T$\cmmnt{, yaitu seluruh ruang yang sedang ditinjau
merupakan himpunan terbuka}. \prooflet{\Prf\
$U(x,1)\subseteq\BbbR^r$ untuk setiap $x\in\BbbR^r$.\ \Qed}

(c) Jika $G$, $H\in\frak T$ maka $G\cap H\in\frak T$\cmmnt{; yaitu irisan
dua himpunan terbuka selalu merupakan himpunan terbuka}.
\prooflet{\Prf\
Misalkan $x\in G\cap H$. Maka terdapat $\delta_1$, $\delta_2>0$ sedemikian
sehingga $U(x,\delta_1)\subseteq G$ dan $U(x,\delta_2)\subseteq H$. Tetapkan
$\delta=\min(\delta_1,\delta_2)>0$; maka

\Centerline{$U(x,\delta)=\{y:\|y-x\|<\min(\delta_1,\delta_2)\}
=U(x,\delta_1)\cap U(x,\delta_2)\subseteq G\cap H$.}

\noindent Karena $x$ sembarang, $G\cap H$ terbuka.\ \Qed}

(d) Jika $\Cal G\subseteq\frak T$, maka

\Centerline{$\bigcup\Cal G=\{x:\,\exists\, G\in\Cal G,\,x\in G\}
\in\frak T$\dvro{.}{;}}

\noindent\cmmnt{yaitu, gabungan sebarang keluarga himpunan terbuka adalah
terbuka.} \prooflet{\Prf\ Misalkan $x\in\bigcup\Cal G$. Maka terdapat suatu
$G\in\Cal G$ sedemikian sehingga $x\in G$. Karena $G\in\frak T$, terdapat
suatu $\delta>0$ sedemikian sehingga


\Centerline{$U(x,\delta)\subseteq G\subseteq\bigcup\Cal G$.}

\noindent Karena $x$ sembarang, $\bigcup\Cal G\in\frak T$.\ \Qed}

\leader{1A2C}{Ketaksamaan Cauchy: Proposisi} Untuk semua
$x$, $y\in\BbbR^r$,
$\|x+y\|\le\|x\|+\|y\|$.

\proof{ Nyatakan $x$ sebagai $(\xi_1,\ldots,\xi_r)$ dan $y$ sebagai
$(\eta_1,\ldots,\eta_r)$;
tetapkan $\alpha=\|x\|$, $\beta=\|y\|$. Maka $\alpha$ dan $\beta$ keduanya
nonnegatif. Jika $\alpha=0$ maka $\sum_{j=1}^r\xi_j^2=0$, sehingga setiap
$\xi_j=0$ dan $x=\tbf{0}$, jadi
$\|x+y\|=\|y\|=\|x\|+\|y\|$; jika $\beta=0$, maka $y=\tbf{0}$ dan
$\|x+y\|=\|x\|=\|x\|+\|y\|$. Jika keduanya positif, tinjau

$$\eqalign{\alpha\beta\|x+y\|^2
&\le\alpha\beta\|x+y\|^2+\|\alpha y-\beta x\|^2\cr
&=\alpha\beta\sum_{j=1}^r(\xi_j+\eta_j)^2
+\sum_{j=1}^r(\alpha\eta_j-\beta\xi_j)^2\cr
&=\sum_{j=1}^r(\alpha\beta\xi_j^2+\alpha\beta\eta_j^2
+\alpha^2\eta_j^2+\beta^2\xi_j^2)\cr
&=\alpha^3\beta+\alpha\beta^3+\alpha^2\beta^2+\beta^2\alpha^2\cr
&=\alpha\beta(\alpha+\beta)^2
=\alpha\beta(\|x\|+\|y\|)^2.\cr}$$

\noindent Dengan membagi kedua ruas oleh $\alpha\beta$ dan mengambil akar
kuadrat, kita memperoleh hasil yang dinyatakan.
}%end of proof of 1A2C

\leader{1A2D}{Akibat} $U(x,\delta)$\cmmnt{, sebagaimana didefinisikan dalam
1A2A,} terbuka untuk setiap $x\in\BbbR^r$ dan $\delta>0$.

\proof{ Jika $y\in U(x,\delta)$, maka $\eta=\delta-\|y-x\|>0$. Sekarang,
jika $z\in U(y,\eta)$,

\Centerline{$\|z-x\|=\|(z-y)+(y-x)\|
\le\|z-y\|+\|y-x\|<\eta+\|y-x\|=\delta$,}

\noindent dan $z\in U(x,\delta)$; dengan demikian
$U(y,\eta)\subseteq U(x,\delta)$. Karena $y$ sembarang,
$U(x,\delta)$ terbuka.
}%end of proof of 1A2D

\leader{1A2E}{Himpunan tertutup: Definisi} Suatu himpunan
$F\subseteq\BbbR^r$ disebut {\bf tertutup} jika $\BbbR^r\setminus F$
terbuka. \cmmnt{({\it Peringatan!} `Sebagian besar' subhimpunan dari
$\BbbR^r$ tidak terbuka maupun tertutup; dua subhimpunan dari $\BbbR^r$,
yaitu $\emptyset$ dan $\BbbR^r$, sekaligus terbuka dan tertutup.) Sejalan
dengan daftar dalam 1A2B, kita mempunyai sifat-sifat berikut dari keluarga
$\Cal F$ yang terdiri atas subhimpunan-subhimpunan tertutup dari $\BbbR^r$.}

\leader{1A2F}{Proposisi} Misalkan $\Cal F$ keluarga subhimpunan-subhimpunan
tertutup dari $\BbbR^r$.

 (a) $\emptyset\in\Cal F$\prooflet{ (karena
$\BbbR^r\in\frak T$)}.

(b) $\BbbR^r\in\Cal F$\prooflet{ (karena $\emptyset\in\frak T$)}.

(c) Jika $E$, $F\in\Cal F$ maka $E\cup F\in\Cal F$\dvro{.}{, karena

\Centerline{$\BbbR^r\setminus(E\cup F)
=(\BbbR^r\setminus E)\cap(\BbbR^r\setminus F)\in\frak T$.}}

(d) Jika $\Cal E\subseteq\Cal F$ merupakan keluarga {\it tak kosong} dari
himpunan-himpunan tertutup, maka

\Centerline{$\bigcap\Cal E
=\{x:x\in F\Forall F\in\Cal E\}\cmmnt{\mskip5mu =\BbbR^r\setminus\bigcupop_{F\in\Cal E}(\BbbR^r\setminus F)}\in\Cal F$.}

\cmmnt{\medskip

\noindent{\bf Catatan} Dalam (d), kita perlu mengasumsikan bahwa
$\Cal E\ne\emptyset$ untuk memastikan bahwa $\bigcap\Cal E\subseteq\BbbR^r$.
}%end of comment

\leader{1A2G}{}\cmmnt{ Sejalan dengan 1A2D, kita memperoleh fakta berikut:

\medskip

\noindent}{\bf Lema} Jika $x\in\BbbR^r$ dan $\delta\ge 0$, maka
$B(x,\delta)=\{y:\|y-x\|\le\delta\}$ tertutup.

\proof{ Tetapkan $G=\BbbR^r\setminus B(x,\delta)$. Jika $y\in G$, maka
$\eta=\|y-x\|-\delta>0$; jika $z\in U(y,\eta)$, maka

\Centerline{$\delta+\eta=\|y-x\|\le\|y-z\|+\|z-x\|<\eta+\|z-x\|$,}

\noindent sehingga $\|z-x\|>\delta$ dan $z\in G$. Jadi,
$U(y,\eta)\subseteq G$. Karena $y$ sembarang, $G$ terbuka dan
$B(x,\delta)$ tertutup.
}%end of proof of 1A2G

\exercises{}

\discrpage
